\section{Introduction}\label{sec:intro}

Weil's criterion relates the Riemann hypothesis to nonnegativity of a
Hermitian form on compactly supported smooth tests
\cite{Weil1952,BOMBIERI-WEILQF-2000}. We use its logarithmic-variable
version, with archimedean, pole and prime contributions. Localized forms
and their variational properties occur in the work of Yoshida,
Bombieri, Connes--Consani and Connes--Consani--Moscovici
\cite{YOSHIDA-HERMITIAN-1992,BOMBIERI-WEILQF-2000,CC-WEILPOS-2020,CCM-ZST-2025}.
The finite CCM matrices are Fourier compressions of a windowed form;
they are not matrices in a prolate eigenfunction basis. Their lowest
eigenvalues give upper bounds for the corresponding full-window
Rayleigh infimum. A finite list of positive upper bounds establishes no
global positivity statement.

We fix Riemann's function $\Phi$, in the convention
$\widehat\Phi=\Xi$, and put $f_0=\Phi/\|\Phi\|_2$
\cite{Riemann1859,CCM-ZST-2025}. The function $f_0$ is smooth and strictly
positive. We specify a completion $\Xf$ on which the Weil form extends
continuously and prove that all translates of $f_0$, and all $f_0*h$ with
$h\in C_c^\infty$, belong to its radical (\Cref{prop:radical}). The cutoff
argument is essential: the explicit formula is not initially being
applied to a compactly supported canonical test.

The substitution $g=f_0s$ yields
\[
 \Q(f_0s)=\int_0^\infty b(t)E_s(t)\,dt
          +\sum_{n\ge2}\frac{\Lam(n)}{\sqrt n}E_s(\log n),
 \qquad
 E_s(t)=\int_{\R}f_0(x)f_0(x+t)|s(x+t)-s(x)|^2dx.
\]
This is a signed difference-form representation, not an assertion of the
positivity or Markov property of a Dirichlet form. The distributional
kernel is already present in Suzuki's screw-function formulas
\cite[\S3.5]{Suzuki2023Screw}\cite[\S2.5]{Suzuki2026}. Ground-state
representations are a classical method, including the nonlocal versions
of Frank and Seiringer \cite{FrankSeiringer2008}. Here we verify the
substitution for a radical vector of a form whose sign is not assumed.
The elementary expression for $b(t)$ follows from the known kernel;
its zero is characterized by the plastic-number equation.
\Cref{cor:DOM} records the resulting three-term positivity inequality.

The main result is \Cref{thm:nominorant}. A fixed nonzero finite sampling
stencil, with a measurable nonnegative weight, cannot minorize
$\Q(f_0s)$ for every compact smooth profile $s$. The proof uses compact
cutoffs of $f_0(\cdot-q)/f_0$, Fatou's lemma, countably many rational
translations and finite translate independence. It applies to an
arbitrary functional with the specified vanishing cutoff values; the
Weil form is a concrete instance. \Cref{cor:noCSS} excludes countable
sums of these finite-stencil energies, not arbitrary sums of squares.
\Cref{ex:control} shows why the obstruction does not imply that the
underlying form has a negative direction.

The quantifier over all compact profiles matters. The known Sonin-space
minorant is imposed on a different, restricted support and moment class
\cite{Connes2026,CC-WEILPOS-2020}; our result does not rule out estimates
on that class. Similarly, \Cref{prop:trade} concerns a fixed margin in
the ambient $\Xf$ norm and says nothing about coercivity in a quotient
norm. The short ledger in \Cref{sec:ledger} keeps these boundaries
explicit. \Cref{sec:windows} reports finite computations and a
conjectural trial-jet expansion, without using them in any proof.
\Cref{sec:disclosure,app:lean} delimit the formalised result and describe
the use of automated tools. No positivity estimate establishing RH is
claimed.

\paragraph{Attribution and scope of the contribution.}
The canonical Fourier representation is classical; the CCM notation
makes its relation to the prolate trial explicit. The same logarithmic
expansion is also recorded in the unrefereed preprint
\cite{Freedman2026}. The signed kernel belongs to the existing
Weil--Suzuki framework, and the ground-state algebra is not a new
general method. The concrete weighted-domain verification and the
global finite-stencil obstruction are the statements developed here.
We have not found the latter in this precise formulation in the sources
cited, but make no exhaustive priority claim. The no-margin proposition
is included as an elementary consequence, not as a separate major
novelty. Related positive-operator and finite-dictionary constructions
\cite{Li2024,Groskin2026} are not inputs asserting positivity of this
Weil form.
